\chapter*{Executive Summary}
\addcontentsline{toc}{chapter}{Executive Summary}
TFE is specified here as a \emph{research instantiation} of DSF-AI. Canonical UF L0--L4 are treated as deterministic, domain-agnostic structural interpreters imported without semantic modification. TFE L5 is the financial governance layer: it is responsible for domain preparation, protected structural memory, multi-timescale interpretation, runtime artifact control, and restrained research-side action formation.

This version corrects a prior deficiency by integrating SES / UF-Core security controls as \emph{normative first-class requirements}, not a mere reference. It also clarifies that TFE computes and governs \emph{structure and geometry} rather than training a direct return-prediction engine, and it separates the scientific thesis from pragmatic implementation compromises by specifying protected operating profiles and explicit L5 conformance profiles. CAP/PRP remain the protected-boundary operating profiles, while CP-0/CP-1/CP-2 distinguish current production, transitional hybrid, and thesis-faithful L5 implementations:
\begin{enumerate}[leftmargin=1.2cm]
    \item \textbf{Conformant Adapter Profile (CAP):} only SES Summary fields cross the protected boundary to standard TFE consumers.
    \item \textbf{Protected Research Profile (PRP):} richer structural payloads are consumed only inside a controlled research boundary where SES protection, chain-of-custody, and auditability remain intact.
\end{enumerate}

The governing thesis for TFE is that, in continuity-rich financial time series, decisions should arise from temporal continuity, field evolution, multi-scale memory, negative space, cross-view consistency, and bounded domain rules, \emph{not} from generic flattened lag tables or unconstrained supervised heuristics. TFE therefore treats any row-first serialized approximation as a pragmatic implementation profile, not as the canonical test of the thesis.

As of v2.5, current production CP-0 is more sharply characterized: decision provenance persistence and full frozen-context parity are green; typed fallback semantics are green; exact-path alignment has been safely adopted and restores L0 exact coverage for a majority of current rows without decision drift; long-side rulebook decomposition and blocked-family contract work are complete; structural recency contract definition is complete but blocked by missing snapshot materialization fields; epoch contract definition is complete but blocked by an unstable numeric regime-code representation, a missing gap-state enum, and an unimplemented sidecar projection join; quote-family activation remains blocked because quote publication is not aligned with the active snapshot and freshness is outside the validation window.

The v2.7 re-architecture resets the system architecture around one deterministic backbone:
\[
\texttt{source} \rightarrow \texttt{normalize} \rightarrow \texttt{assess candidate} \rightarrow \texttt{publish immutable bundle} \rightarrow \texttt{serve active bundle}
\]
This revision treats TFE as a financial advisor platform and establishes the Source Plane, Assessment and Publication Plane, Serving Plane, Enrichment Plane, Refresh State Machine, Product Deployment Plane, and Migration/Cutover as normative architecture.

This consolidated specification:
\begin{enumerate}[leftmargin=1.2cm]
    \item imports canonical UF L0--L4 and their financial adapter boundary;
    \item normatively integrates SES/UF-Core security, envelope protection, and threat responses;
    \item defines a deterministic multi-view, multi-core, event-driven L5 with shared long/mid/short memory;
    \item makes explicit how TFE computes geometric structure, event-tape dynamics, and bounded decision support differently from standard predictive analytics and supervised ML;
    \item defines the deterministic source-assess-publish-serve architecture with immutable publication bundles and phase-ledger state machine;
    \item defines finance-specific rules, admissible heuristics, strategy classes, epoch mosaics, sector/sphere coupling, runtime snapshots, persistence, auditability, and validation;
    \item defines conformance, benchmark, and release rules, including a commercial kill rule; and
    \item states applicability boundaries and non-goals, including the distinction between continuity-rich time-series problems and symbolic/combinatorial tasks better served by other methods.
\end{enumerate}

\chapter{Document Control, Purpose, and Scope}
\section{Purpose}
\meta{State why TFE exists and what this specification governs.}{UF v1.4.0 kernel definition, UF-Core/SES security specification, financial domain constraints, DSF-AI thesis, governance requirements.}{Controlled purpose statement and scope boundary.}

The purpose of TFE is to instantiate DSF-AI in a financial research setting while preserving the canonical role of UF L0--L4 as structural interpreters. TFE does not redefine UF. It constrains how UF outputs are invoked, protected, persisted, interpreted over time, validated, and transformed into governed research signals.

The purpose of the v2.7 architecture layer is to define the target deterministic architecture for TFE as a commercially operable advisor system with auditable publication semantics, strict stage boundaries, and exact runtime state visibility.

\section{Scope}
\meta{Define the boundaries of the present specification.}{Financial time series, UF kernel outputs, protected runtime artifacts, validation controls, approved rules.}{In-scope and out-of-scope statements.}

In scope:
\begin{itemize}
    \item financial-domain input preparation and deterministic invocation of canonical UF L0--L4;
    \item SES / UF-Core protection over UF-SP and derivative protected research representations;
    \item TFE L5 state, governance, action gating, multi-timescale memory, runtime snapshot semantics, persistence, validation, and auditability;
    \item controlled use of richer structural payloads only inside the protected research profile;
    \item source acquisition and normalization contracts;
    \item candidate build, assessment, publication, activation, and serving contracts;
    \item refresh and phase-ledger state machines;
    \item deployment/release plane separation;
    \item deterministic identifiers, invariants, and failure semantics;
    \item validation, traceability, and conformance rules.
\end{itemize}

Out of scope:
\begin{itemize}
    \item autonomous trading or brokerage execution;
    \item modification of canonical UF kernel semantics;
    \item unsupported benchmark claims or commercial statements without artifact-backed evidence;
    \item symbolic/combinatorial prediction domains in which the sequential structural assumptions of DSF-AI do not hold without a separate domain engine;
    \item marketing claims;
    \item aspirational functionality without a declared contract;
    \item uncontrolled heuristic behavior;
    \item undocumented cross-boundary shortcuts;
    \item microsecond trading-latency optimization as the principal architecture target.
\end{itemize}

\section{Roles and Responsibilities}
\meta{Define the responsibility classes required by good documentation practice and controlled execution.}{Project governance structure and quality responsibilities.}{Role definitions and control obligations.}

At minimum, the following roles shall exist:
\begin{itemize}
    \item \textbf{Specification Owner:} accountable for scope, approvals, and controlled revision of this specification.
    \item \textbf{Kernel Steward:} responsible for UF version compatibility, import correctness, and non-modification of canonical semantics.
    \item \textbf{Security Steward:} responsible for SES integration, envelope controls, keying assumptions, threat response implementation, and chain-of-custody security events.
    \item \textbf{L5 Steward:} responsible for TFE governance logic, multi-scale memory design, and domain-rule correctness.
    \item \textbf{Validator:} responsible for conformance, reproducibility, benchmark execution, and release disposition.
    \item \textbf{Records Steward:} responsible for audit trails, retention, and controlled records.
\end{itemize}

\section{TFE as a Research Instantiation of DSF-AI}
\meta{Locate TFE inside the broader DSF-AI program.}{DSF-AI thesis, canonical UF kernel, domain-specific governance requirements.}{A precise identity for TFE.}

TFE is a domain-specific research instantiation of DSF-AI. In TFE:
\begin{enumerate}[leftmargin=1.2cm]
    \item UF L0--L4 provide deterministic structural interpretation of an ordered financial window.
    \item TFE L5 maintains governed memory over sequences of structural events and runtime metadata.
    \item Approved domain rules determine when structural states are admissible for external research interpretation.
    \item Outputs are qualitative research signals with abstention and fail-safe behavior, not autonomous execution directives.
\end{enumerate}

\section{Applicability Boundary and Complementarity}
\meta{State where TFE is expected to add value and where complementary methods remain appropriate.}{Domain characteristics, sequence continuity, symbolic density, validation goals.}{Explicit applicability boundary.}

TFE is designed for continuity-rich sequential domains in which structural stability, robustness, phase transition, collapse, event ordering, temporal continuity, negative space, and coupled fast/medium/slow memory are meaningful.

A useful shorthand is that DSF-AI is strongest for the ``1/2/3'' aspects of a domain --- continuous ordered evolution, structural geometry, and phase-sensitive field dynamics --- whereas conventional ML or symbolic/combinatorial engines remain valuable for the ``A/B/C'' aspects --- discrete class labeling, large symbolic search, or rule-dense combinatorics. TFE is therefore specified as a complementary research system, not a blanket replacement for every modeling paradigm.
